\documentclass[11pt,a4paper]{article}
\usepackage[T1]{fontenc}
\usepackage[left=2cm, right=2cm, top=2cm, bottom=2cm]{geometry}
\usepackage{graphicx}
\usepackage{mathtools}
\usepackage{amssymb}
\usepackage{amsthm}
\usepackage{thmtools}
\usepackage{xcolor}
\usepackage{nameref}
\usepackage[colorlinks=true, linkcolor=black, urlcolor=black, citecolor=black]{hyperref}

\usepackage{physics}
\usepackage{amsmath}
\usepackage{tikz}
\usepackage{mathdots}
\usepackage{yhmath}
\usepackage{cancel}
\usepackage{color}
\usepackage{siunitx}
\usepackage{array}
\usepackage{multirow}
\usepackage{amssymb}
\usepackage{gensymb}
\usepackage{tabularx}
\usepackage{extarrows}
\usepackage{booktabs}
\usetikzlibrary{fadings}
\usetikzlibrary{patterns}
\usetikzlibrary{shadows.blur}
\usetikzlibrary{shapes}

\usepackage{placeins}




\theoremstyle{definition}
\newtheorem{defn}{Definition}


\title{A Second Chance For Random Varaibles}
\author{Ali Fele Paranj}

\newtheoremstyle{Example}
{\topsep}   % Space above
{\topsep}   % Space below
{}  % Body font
{}          % Indent amount
{\bfseries} % Theorem head font (bold)
{.}         % Punctuation after theorem head
{ }         % Space after theorem head
{}          % Theorem head spec


\theoremstyle{Example}
\newtheorem{example}{Example}



\newcommand{\PP}{\mathbb{P}}
\newcommand{\set}[1]{\{#1\}}
\newcommand{\N}{\mathbb{N}}




\begin{document}
	\maketitle
	\begin{abstract}
		The notions of random variables, random numbers, sampling a distribution, and etc are highly abused in different fields with numerous mix-and-match in notations and definitions. In this paper I will clarify some of these notions such that one can distinguish between the idea of a random number, a random variable, the meaning of sampling a distribution.
	\end{abstract}
	
	\section{Introduction}
	What is a random variable? If you look up in a ``proper''\footnote{i.e. written by a ``proper'' (whatever that means) mathematician.} probability textbook, you would encounter something like the following definition.
	
	\begin{defn}
		Given a probability space $ (\Omega,\mathcal{F},\mathbb{P}) $, a random variable $ X $ is a measurable function from $ (\Omega, \mathcal{F}) $ to a measurable space $ (S,\mathcal{S}) $. 
	\end{defn}
	
	However, when you want to work with probability theory in a computer or an applied mathematics class, then you might encounter that the instructor use the notion of the random variable interchangeably with the notion of random number generated by some pseudo random number generator algorithm. And they might use some ideas like ``realization of a random variable'', in which a random variable magically turns into some random numbers! Or they might talk about sampling a distribution (or in some classes I've heard about sampling a random variable!). But what does it really mean? What does it mean to sample a distribution? What does it mean to sample via a Markov chain? How these ``applied'' ideas of random numbers, realizations, etc relates to random variables, that are just functions? Neither of these questions are answered properly, and often leaves the students with many ambiguous answers.
	
	
	
	
	
	\section{Construction of Appropriate Sample Space}
	Consider the set $ S = \set{0,1} $. Let $ \Omega = S^\N $ be the product space with $ X_n:\Omega \to S $ as the projection maps, which are Bernoulli random variables with parameter 1/2. I.e.
	\[ \PP(X_n^{-1}(0)) = \PP(X_n^{-1}(0)) = 1/2. \]
	Intuitively speaking $ \Omega $ is the set of all sequences of zeros and ones and $ X_n $ is the $ n^\text{th} $ element of the sequence, and exactly half the sequences has in their $ n^\text{th} $ position value $ 0 $ and the other half has value $ 1 $. One might find this set $ \Omega $ and the projection maps $ (X_n)_n $ very simple, but this will be the central object of our discussion.
	
	We can define other sequence of functions (i.e. random variables) on $ \Omega $ that would partition $ \Omega $ with different weights. Each of the random variables in $ (X_n)_{n\in\N} $ partition $ \Omega $ into two ``equal'' partition $ X_n^{-1}(0) $ and $ X_n^{-1}(1) $. So we want to construct new sequence of random variables that each of them in the sequence partitions $ \Omega $ into different ``sized'' pieces.
	
	
	\subsection{Uniform Random Variables}
	
	
	\textbf{Random variable on $ \set{0,\cdots,2^{k}-1} $:} Assume we want to construct random variables $ (U_n)_{n\in N} $ defined on $ \Omega $ such that has uniform distribution over $ \set{0,1,\cdots,2^k-1} $ (i.e. its CDF jumps at every integer $ 0,\cdots,2^k-1 $ by $ 1/2^k $ being right continuous at each jump). In other words we want
	\[ \PP(\set{U_n = i}) = 1/2^k \qquad \text{for $ i\in\set{0,1,\cdots,2^k-1} $}. \]
	
	We claim that the following random variable does the job. Define $ U_n:\Omega\to\set{0,1,\cdots,2^k-1} $ as follows:
	\[U_n(\omega) = \overline{X_{nk}(\omega) X_{nk+1}(\omega) X_{nk+2}(\omega)\cdots X_{nk+k-1}(\omega)}, \]
	where the over bar indicates that we mean the binary representation of a number whose bits are given by the random variables. For instance, let $ \omega = (0,1,1,0,1,0,1,0,0,0,0,1,0,0,1,1,0,1,0,0,\cdots) $. Then 
	\[ U_0(\omega) = \overline{011}=3, \qquad U_1(\omega) =\overline{010}=2, \qquad U_3(\omega) = \overline{100}=4, \qquad U_4(\omega) = \overline{001} = 1, \qquad \cdots. \]
	
	So we claim that for each $ n\in\N $ the random variable $ U_n $ is a uniform random variable over $ \set{0,\cdots,2^k-1} $. To see this observe that for $ i = \overline{a_0a_1\cdots a_k} $
	\[ \PP(B_n(\omega) = i) = \PP(X_0=a_0,X_1=a_1,\cdots,X_k=a_k) = 1/2^k. \]
	
	
	\textbf{Random variable on $ \set{0,\frac{1}{2^k},\frac{2}{2^k},\cdots,\frac{2^k-1}{2^k}}\subset[0,1] $:} In this case we want the random variable to assume values in $ S =  \set{0,\frac{1}{2^k},\frac{2}{2^k},\cdots,\frac{2^k-1}{2^k}}\subset[0,1]$. The binary expansion of $ i/2^k \in S $ is the same as $ i $ but shifted $ k $ positions to left. For instance, if the binary expansion of $ i $ is of the form $\overline{ a_0\cdots a_n} $, then they binary expansion of $ i/2^k $ is $ \overline{0.a_0\cdots a_n} $. So this time, we define the random variable $ U_n $ to be
	\[U_n(\omega) = \overline{0.X_{nk}(\omega) X_{nk+1}(\omega) X_{nk+2}(\omega)\cdots X_{nk+k-1}(\omega)}. \]
	As similar idea as above can be used to show that $ \PP(U_n = i/2^k) = 1/2^k  $ for all $ i/2^k \in S $.
	

	
	\subsection{Arbitrary Random Variables}
	There is no one-fits-all method to construct a sequence of random variable on $ \Omega $ that each of them has arbitrary distribution. Constructing such sequence of random variables if also often called as \textbf{sampling a distribution}. So there are different methods to sample a distribution. One well-known and simple method is inverse transform sampling.
	
	
	\textbf{Inverse transform sampling:} Assume we want to construct a sequence of  random variables $ (X_n)_{n\in\N} $ where for each $ n\in \N $ $ X_n:\Omega \to \set{0,1,2,3} $ where
	\[ \PP(X_n^{-1}(0)) = \frac{6}{10}, \qquad \PP(X_n^{-1}(1)) = \frac{1}{10},\qquad \PP(X_n^{-1}(2)) = \frac{2}{10},\qquad \PP(X_n^{-1}(3)) = \frac{1}{10}, \]
	a.k.a we want to draw samples from the distribution $ \mu:\set{0,1,2,3}\to[0,1] $ where
	\[ \mu(0) = \frac{6}{10}, \qquad \mu(1) = \frac{1}{10}, \qquad \mu(2)=\frac{2}{10}, \qquad \mu(3) = \frac{1}{10}. \]
	a.k.a. we want to construct $ (X_n)_{n\in\N} $ such that for all $ n $ $ X_n $ has the following CDF
	\[ F(x) = \begin{cases}
		0 \qquad x<0,\\
		\frac{6}{10} \qquad 0\leq x < 1,\\
		\frac{7}{10} \qquad 1\leq x < 2,\\
		\frac{9}{10} \qquad 2\leq x < 3,\\
		1 \qquad 3\leq x,\\
	\end{cases} \]
	The idea is to calculate the ``inverse'' of the function $ F $ and define $ X_n $ as $ X_n = F^{-1} \circ U_n $. Note that $ F(x) $ might fail to have an inverse, the discontinuities of $ F $ turn into constant values in $ F^{-1} $ and the points of constant value in $ F $ turn into discontinuity if $ F^{-1} $. This is best described by a picture.
	\begin{figure}[h!]
	\centering
	
	
	
	
	\tikzset{every picture/.style={line width=0.75pt}} %set default line width to 0.75pt        
	
	\begin{tikzpicture}[x=0.75pt,y=0.75pt,yscale=-1,xscale=1]
		%uncomment if require: \path (0,300); %set diagram left start at 0, and has height of 300
		
		%Straight Lines [id:da966046102014967] 
		\draw    (50,200) -- (258,200) ;
		\draw [shift={(260,200)}, rotate = 180] [color={rgb, 255:red, 0; green, 0; blue, 0 }  ][line width=0.75]    (10.93,-3.29) .. controls (6.95,-1.4) and (3.31,-0.3) .. (0,0) .. controls (3.31,0.3) and (6.95,1.4) .. (10.93,3.29)   ;
		%Straight Lines [id:da6895762249704231] 
		\draw    (100,210) -- (100,72) ;
		\draw [shift={(100,70)}, rotate = 90] [color={rgb, 255:red, 0; green, 0; blue, 0 }  ][line width=0.75]    (10.93,-3.29) .. controls (6.95,-1.4) and (3.31,-0.3) .. (0,0) .. controls (3.31,0.3) and (6.95,1.4) .. (10.93,3.29)   ;
		%Straight Lines [id:da7623069951769825] 
		\draw [color={rgb, 255:red, 74; green, 144; blue, 226 }  ,draw opacity=1 ][line width=1.5]    (50,200) -- (98.39,200) ;
		\draw [shift={(100,200)}, rotate = 0] [color={rgb, 255:red, 74; green, 144; blue, 226 }  ,draw opacity=1 ][line width=1.5]      (0, 0) circle [x radius= 2.61, y radius= 2.61]   ;
		%Straight Lines [id:da5999081723558621] 
		\draw [color={rgb, 255:red, 74; green, 144; blue, 226 }  ,draw opacity=1 ][line width=1.5]    (100,140) -- (128.39,140) ;
		\draw [shift={(130,140)}, rotate = 0] [color={rgb, 255:red, 74; green, 144; blue, 226 }  ,draw opacity=1 ][line width=1.5]      (0, 0) circle [x radius= 2.61, y radius= 2.61]   ;
		\draw [shift={(100,140)}, rotate = 0] [color={rgb, 255:red, 74; green, 144; blue, 226 }  ,draw opacity=1 ][fill={rgb, 255:red, 74; green, 144; blue, 226 }  ,fill opacity=1 ][line width=1.5]      (0, 0) circle [x radius= 2.61, y radius= 2.61]   ;
		%Straight Lines [id:da13557795475547252] 
		\draw [color={rgb, 255:red, 74; green, 144; blue, 226 }  ,draw opacity=1 ][line width=1.5]    (130,130) -- (158.39,130) ;
		\draw [shift={(160,130)}, rotate = 0] [color={rgb, 255:red, 74; green, 144; blue, 226 }  ,draw opacity=1 ][line width=1.5]      (0, 0) circle [x radius= 2.61, y radius= 2.61]   ;
		\draw [shift={(130,130)}, rotate = 0] [color={rgb, 255:red, 74; green, 144; blue, 226 }  ,draw opacity=1 ][fill={rgb, 255:red, 74; green, 144; blue, 226 }  ,fill opacity=1 ][line width=1.5]      (0, 0) circle [x radius= 2.61, y radius= 2.61]   ;
		%Straight Lines [id:da8593539343859463] 
		\draw [color={rgb, 255:red, 74; green, 144; blue, 226 }  ,draw opacity=1 ][line width=1.5]    (160,110) -- (188.39,110) ;
		\draw [shift={(190,110)}, rotate = 0] [color={rgb, 255:red, 74; green, 144; blue, 226 }  ,draw opacity=1 ][line width=1.5]      (0, 0) circle [x radius= 2.61, y radius= 2.61]   ;
		\draw [shift={(160,110)}, rotate = 0] [color={rgb, 255:red, 74; green, 144; blue, 226 }  ,draw opacity=1 ][fill={rgb, 255:red, 74; green, 144; blue, 226 }  ,fill opacity=1 ][line width=1.5]      (0, 0) circle [x radius= 2.61, y radius= 2.61]   ;
		%Straight Lines [id:da18454051252823822] 
		\draw [color={rgb, 255:red, 74; green, 144; blue, 226 }  ,draw opacity=1 ][line width=1.5]    (190,100) -- (240,100) ;
		\draw [shift={(190,100)}, rotate = 0] [color={rgb, 255:red, 74; green, 144; blue, 226 }  ,draw opacity=1 ][fill={rgb, 255:red, 74; green, 144; blue, 226 }  ,fill opacity=1 ][line width=1.5]      (0, 0) circle [x radius= 2.61, y radius= 2.61]   ;
		%Straight Lines [id:da8429762985239627] 
		\draw  [dash pattern={on 0.84pt off 2.51pt}]  (130,130) -- (130,140) ;
		%Straight Lines [id:da8056681225243965] 
		\draw  [dash pattern={on 0.84pt off 2.51pt}]  (160,110) -- (160,130) ;
		%Straight Lines [id:da7911133706970217] 
		\draw  [dash pattern={on 0.84pt off 2.51pt}]  (190,100) -- (190,110) ;
		%Straight Lines [id:da9129197315594763] 
		\draw    (330,200) -- (538,200) ;
		\draw [shift={(540,200)}, rotate = 180] [color={rgb, 255:red, 0; green, 0; blue, 0 }  ][line width=0.75]    (10.93,-3.29) .. controls (6.95,-1.4) and (3.31,-0.3) .. (0,0) .. controls (3.31,0.3) and (6.95,1.4) .. (10.93,3.29)   ;
		%Straight Lines [id:da14882188471018387] 
		\draw    (350,210) -- (350,72) ;
		\draw [shift={(350,70)}, rotate = 90] [color={rgb, 255:red, 0; green, 0; blue, 0 }  ][line width=0.75]    (10.93,-3.29) .. controls (6.95,-1.4) and (3.31,-0.3) .. (0,0) .. controls (3.31,0.3) and (6.95,1.4) .. (10.93,3.29)   ;
		%Straight Lines [id:da6759430187821265] 
		\draw [color={rgb, 255:red, 0; green, 0; blue, 0 }  ,draw opacity=1 ][line width=1.5]    (410,200) -- (350,200) ;
		%Straight Lines [id:da5198928816022348] 
		\draw [color={rgb, 255:red, 0; green, 0; blue, 0 }  ,draw opacity=1 ][line width=1.5]    (420,170) -- (410,170) ;
		%Straight Lines [id:da26153884398935556] 
		\draw [color={rgb, 255:red, 0; green, 0; blue, 0 }  ,draw opacity=1 ][line width=1.5]    (440,140) -- (420,140) ;
		%Straight Lines [id:da12929802874318563] 
		\draw [color={rgb, 255:red, 0; green, 0; blue, 0 }  ,draw opacity=1 ][line width=1.5]    (450,110) -- (440,110) ;
		%Straight Lines [id:da32612822871855307] 
		\draw [color={rgb, 255:red, 74; green, 144; blue, 226 }  ,draw opacity=1 ] [dash pattern={on 0.84pt off 2.51pt}]  (410,170) -- (410,200) ;
		%Straight Lines [id:da8379258814317788] 
		\draw [color={rgb, 255:red, 74; green, 144; blue, 226 }  ,draw opacity=1 ] [dash pattern={on 0.84pt off 2.51pt}]  (420,140) -- (420,170) ;
		%Straight Lines [id:da7260941442334858] 
		\draw [color={rgb, 255:red, 74; green, 144; blue, 226 }  ,draw opacity=1 ] [dash pattern={on 0.84pt off 2.51pt}]  (440,110) -- (440,140) ;
		%Straight Lines [id:da6344873542374194] 
		\draw    (130,195) -- (130,205) ;
		%Straight Lines [id:da9481093762792746] 
		\draw    (160,195) -- (160,205) ;
		%Straight Lines [id:da3399422030709628] 
		\draw    (190,195) -- (190,205) ;
		%Straight Lines [id:da5025910664038975] 
		\draw    (220,195) -- (220,205) ;
		%Straight Lines [id:da673957541783882] 
		\draw    (355,110) -- (345,110) ;
		%Straight Lines [id:da5023584214975627] 
		\draw    (105,100) -- (95,100) ;
		%Straight Lines [id:da38251647502892805] 
		\draw    (355,170) -- (345,170) ;
		%Straight Lines [id:da638486370407335] 
		\draw    (355,140) -- (345,140) ;
		%Straight Lines [id:da23901540996445592] 
		\draw    (410,195) -- (410,205) ;
		%Straight Lines [id:da9038102040428887] 
		\draw    (420,195) -- (420,205) ;
		%Straight Lines [id:da39386684006506933] 
		\draw    (440,195) -- (440,205) ;
		%Straight Lines [id:da11224396769802392] 
		\draw    (450,195) -- (450,205) ;
		
		% Text Node
		\draw (87,98.4) node [anchor=north west][inner sep=0.75pt]  [font=\tiny]  {$1$};
		% Text Node
		\draw (127,207.4) node [anchor=north west][inner sep=0.75pt]  [font=\tiny]  {$1$};
		% Text Node
		\draw (157,207.4) node [anchor=north west][inner sep=0.75pt]  [font=\tiny]  {$2$};
		% Text Node
		\draw (187,207.4) node [anchor=north west][inner sep=0.75pt]  [font=\tiny]  {$3$};
		% Text Node
		\draw (217,207.4) node [anchor=north west][inner sep=0.75pt]  [font=\tiny]  {$4$};
		% Text Node
		\draw (93,207.4) node [anchor=north west][inner sep=0.75pt]  [font=\tiny]  {$0$};
		% Text Node
		\draw (337,167.4) node [anchor=north west][inner sep=0.75pt]  [font=\tiny]  {$1$};
		% Text Node
		\draw (337,137.4) node [anchor=north west][inner sep=0.75pt]  [font=\tiny]  {$2$};
		% Text Node
		\draw (337,107.4) node [anchor=north west][inner sep=0.75pt]  [font=\tiny]  {$3$};
		% Text Node
		\draw (400,207.4) node [anchor=north west][inner sep=0.75pt]  [font=\tiny]  {$0.6$};
		% Text Node
		\draw (414,207.4) node [anchor=north west][inner sep=0.75pt]  [font=\tiny]  {$0.7$};
		% Text Node
		\draw (431,207.4) node [anchor=north west][inner sep=0.75pt]  [font=\tiny]  {$0.9$};
		% Text Node
		\draw (448,207.4) node [anchor=north west][inner sep=0.75pt]  [font=\tiny]  {$1$};
		
		
	\end{tikzpicture}
\end{figure}
	Observe that since $ X_n = F^{-1}\circ U_n $, and $ F^{-1} $ is an step function, we only need enough bits in $ U_n $ to determine which bin it falls into.
	\FloatBarrier
	
	
	\textbf{Sampling from exponential distribution:} The CDF for an exponential distribution is 
	\[ F(x) = \begin{cases}
		0 \qquad  x< 0, \\
		1-e^{-\lambda x} \qquad 0\leq x.
	\end{cases} \]
	This function is invertible and the inverse is
	\[ F^{-1}(x) = -\frac{1}{\lambda}\ln(1-x) \]
	And by defining $ X_n = F^{-1}(U_n) $, we have a sequence of random variables that for each $ n\in\N $ $ X_n $ exponential distribution.
	\[ F_{X_n}(x) = \PP(X_n \leq x) = \PP(F(X_n) \leq F(x)) = \PP(U_n \leq F(x)) = F(x) \]
	
	
	
	\section{On Computers}
	In short, in order to sample from a distribution, we need to construct a sequence of random variables on $ \Omega = \set{0,1}^\N $ such that $ X_n $ has the required distribution. For instance, $ X_n $ can look at the the bits in from position $ nf(n) $ to $ (n+1)f(n) $, and report appropriate number. The different sampling methods help in designing such $ X_n $ on $ \Omega $

\end{document}

